% ============================================================================
% Discussion
%
% Purpose
% -------
% Explain what the results mean.
%
% Reader should leave knowing:
%   • why the findings matter
%   • how they compare with previous research
%   • educational implications
%   • strengths and limitations
%
% Do NOT:
%   • repeat the Results
%   • introduce new analyses
%
% Transition:
% Lead naturally to the Conclusion.
% ============================================================================


\section{Discussion}
\label{sec:discussion}

The results of this study provide new insights into the factors associated with successful transition into first-stage university mathematics. 
By combining reproducible analytical methods with interpretable machine learning models, the analysis extends beyond prediction alone to examine 
the relative importance and consistency of the variables associated with student success. This discussion considers these findings in the context 
of previous research, reflects on their educational significance, and identifies implications for both institutional practice and future research.

Secondary mathematics preparation emerged as the strongest and most influential predictor group across the tuned models, with a mean feature importance 
of \SecondaryMathImportance. Composite risk (\CompositeRiskImportance) and preparatory mathematics (\PreparatoryMathImportance) were the next most influential 
predictor groups. Collectively, these three groups accounted for \TopThreeImportance\% of the average feature importance, indicating that measures of prior 
mathematical preparation dominated the predictive performance of the evaluated models.
Furthermore, the consistency of the feature rankings presented in Section~\ref{sec:feature_rankings} showed that secondary mathematics preparation remained 
influential across all \TunedModelCountWord\ tuned models, whereas the composite risk measure contributed strongly but appeared in only two models. 
Together, these findings provide robust evidence that students' secondary mathematical preparation was the strongest and most consistent predictor of success in this cohort.

These findings suggest that successful transition into first-stage university mathematics is influenced primarily by students' existing mathematical 
preparedness rather than by demographic or administrative characteristics alone. The consistency of secondary mathematics preparation across multiple 
modelling approaches reinforces the importance of prior mathematical knowledge as a foundation for university success. At the same time, the contribution 
of preparatory mathematics and the composite risk measure indicates that institutional interventions can play a meaningful role in identifying and supporting 
students whose prior preparation may place them at greater academic risk.

\subsection*{Relationship to Previous Research}

Previous research has consistently demonstrated that students' prior mathematical preparation influences their transition to university 
mathematics~\cite{Hourigan2007underpreparedness,Wood2001,England2021transition}. This includes previous work undertaken by the author \cite{Stanley2025}. 
The findings suggest that institutions should continue to place strong emphasis on identifying students' mathematical preparedness before or at the 
commencement of university study~\cite{Reddy02122025}. Rather than contradicting existing evidence, the results strengthen current understanding by demonstrating that these relationships remain evident 
across multiple analytical approaches \cite{Hourigan2007underpreparedness,Wood2001,England2021transition}.

\subsection*{Contribution of the Present Study}

The principal contribution of this study is not the identification of prior mathematical preparation as an important predictor of student success, since this has been 
consistently reported in previous transition research. Rather, the contribution lies in the development and application of a transparent, reproducible analytical workflow 
that integrates statistical analysis, interpretable machine learning, independent verification, and automated research reporting. This framework enables educational 
researchers to investigate transition outcomes using methods that are both explainable and reproducible, while maintaining a clear connection between quantitative 
evidence and educational interpretation.

\subsection*{Educational Implications}

The findings suggest that institutions should continue to place strong emphasis on identifying students' mathematical preparedness before or at the commencement of 
university study. Rather than viewing prior preparation as a fixed limitation, the results support the use of early diagnostic measures and targeted preparatory programs 
to identify students who may benefit from additional support. The prominence of preparatory mathematics and the composite risk measure indicates that institutional 
interventions can complement students' existing knowledge by providing structured pathways that improve their likelihood of success.

\subsection*{Limitations and Future Research}

The present study should be interpreted within the context of several limitations. The analysis was conducted using data from a single first-stage university mathematics 
subject and therefore reflects the characteristics of one institutional context rather than the broader higher education sector. Although the analytical workflow was 
intentionally designed to be reproducible and transferable, the specific predictive relationships identified in this study may vary across institutions, disciplines, and 
student cohorts. Future research should apply the framework to additional subjects, longitudinal datasets, and diverse institutional settings to evaluate its broader 
generalisability while continuing to refine interpretable approaches for educational learning analytics.

